\documentclass[11pt]{article}
% Shared preamble for the hanzo-kernel Paper 07 series.
% One style, one vocabulary, inputted by every paper so the series reads as one voice.
\usepackage[margin=1in]{geometry}
\usepackage{booktabs}
\usepackage{amsmath}
\usepackage{amssymb}
\usepackage{graphicx}
\usepackage{xcolor}
\usepackage{hyperref}
\hypersetup{colorlinks=true, linkcolor=blue, urlcolor=blue, citecolor=blue}
\usepackage{enumitem}
\usepackage{listings}
\usepackage{array}

\definecolor{kw}{rgb}{0.13,0.29,0.53}
\definecolor{cmt}{rgb}{0.40,0.40,0.40}
\definecolor{str}{rgb}{0.16,0.50,0.16}
\lstset{
  basicstyle=\ttfamily\footnotesize,
  keywordstyle=\color{kw}\bfseries,
  commentstyle=\color{cmt}\itshape,
  stringstyle=\color{str},
  showstringspaces=false,
  breaklines=true,
  columns=fullflexible,
  frame=single,
  framesep=4pt,
}
\lstdefinelanguage{Rust}{
  morekeywords={fn,let,mut,pub,struct,impl,trait,enum,match,if,else,for,while,loop,return,use,mod,crate,self,super,async,await,where,type,const,static,unsafe,ref,move,dyn,as,true,false},
  morecomment=[l]{//}, morecomment=[s]{/*}{*/}, morestring=[b]", sensitive=true,
}

\newcommand{\x}{\ensuremath{\times}}
\newcommand{\dpa}{\texttt{dp4a}}
\newcommand{\hk}{\texttt{hanzo-kernel}}
\newcommand{\code}[1]{\texttt{#1}}

% Absorb minor prose overfulls without rivers; keep line-breaking sane.
\tolerance=800
\emergencystretch=3em


\usepackage{amsthm}
\newtheorem{proposition}{Proposition}
\newtheorem{lemma}{Lemma}
\newtheorem{definition}{Definition}
\newtheorem{remark}{Remark}

\newcommand{\E}{\mathbb{E}}
\newcommand{\R}{\mathbb{R}}
\newcommand{\avg}[1]{\langle #1 \rangle}

\title{\textbf{Evolutionary Orchestration of Agentic Fleets:\\ Population Search, Recombination, and the Explore--Exploit Anneal}}
\author{Hanzo AI --- ML Systems}
\date{hanzo-agentic Paper 02 --- the \emph{Continually Improving Fleets} series}

\begin{document}
\maketitle

\begin{abstract}
An agentic fleet is a population of candidate behaviours acting on a shared solution space. We argue that
the natural control law for such a fleet is not command but \emph{evolution}: selection, recombination, and
mutation applied to the population, generation over generation. The claim is made precise by placing three
folk mechanisms on one object --- a probability distribution $p_t$ over solution space --- and showing they
are the same dynamics at three timescales. A mean-field game (companion paper,
\emph{Mean-Field Alignment}) is the fast inner loop: each agent best-responds to the population mean and the
fleet equilibrates. Ferromagnetic criticality is the thermostat: it sets how aligned the population is.
A genetic algorithm (GA) is the slow outer loop: it moves $p_t$ itself uphill on fitness. We show a GA is a
fitness-driven Fokker--Planck flow --- selection is replicator drift, mutation is diffusion, crossover is a
mixing operator --- and that when fitness is frequency-dependent, replicator drift \emph{is} the mean-field
coupling of the game. The one operator the other two lack is crossover, and it is the source of the speedup:
by the schema theorem it recombines above-average building blocks, so on a \emph{separable} landscape it
assembles the global optimum in time polynomial in the number of blocks where mutation-only takes
exponential time (Proposition~\ref{prop:speedup}); the honest boundary is deception, where crossover is
disruptive. Premature convergence is exactly the ferromagnetic ordered phase (groupthink), and the cure is
annealing, not quenching. We then ground the framework twice against real Hanzo systems: the GRPO rollout
\emph{group} of \cite{hanzo-native-training} is a GA population and its group-relative advantage is selection;
AnyMoE experts are building blocks and expert recombination is crossover; the island model of
\cite{hanzo-intrinsic-islands} is a coarse-grained GA. In the agentic harness the population is a set of
solution \emph{variants}, fitness is the adversarial critic together with a green build, crossover is a
synthesis agent that grafts the fittest module from each parent, and generations are the loop-until-dry
rounds. We keep the accounting honest: this is a control lens, not a literal energy function, but the
building-block argument is a genuine reason recombination accelerates the decomposable problems --- files,
modules, tools --- that engineering fleets actually solve.
\end{abstract}

\section{The pain: a fleet that averages when it should assemble}
\label{sec:pain}

Run $N$ agents at a hard problem and the two obvious control laws both disappoint. Make them independent and
they duplicate each other's dead ends; nobody's partial success helps anybody else, and the fleet is a set of
lonely random walks. Make them agree --- align them to a running consensus --- and they converge, but only to
the \emph{mean} of what they already believed; consensus averages, it does not invent. Neither law does the
thing a research group does when it is working well: take the best idea from one line of attack, the best idea
from another, and \emph{combine} them into something better than either.

That combining operation has a name. It is crossover, and it is the defining move of a genetic algorithm.
This paper's thesis is that an agentic fleet should be run as a population under evolutionary search, that
crossover is precisely the operator the ``independent'' and ``aligned'' regimes lack, and that it is fast for
a concrete, checkable reason: the problems a software fleet solves are \emph{decomposable} into modules, and
crossover is the operator that exploits decomposition. We make each of those claims formal, then show the
same three knobs already live inside Hanzo's model-training loop and its agent harness --- so ``continually
improving fleets'' is not an aspiration bolted on the side but the loop the substrate already runs, named.

\section{One population, three timescales}
\label{sec:frame}

Fix a solution space $\mathcal{S}$ (policies, programs, tool implementations --- whatever a unit of the fleet
emits) and let the state of the fleet be a distribution $p_t \in \Delta(\mathcal{S})$. Everything below is a
map $p_t \mapsto p_{t+1}$. Three such maps act at separated timescales, and separation of timescales is what
lets them compose without fighting.

\begin{definition}[The three loops]
\label{def:loops}
\emph{Inner (fast) --- the game.} Holding $p_t$ fixed, each agent chooses a best response to the induced mean
field; iterating to the fixed point is a mean-field-game (MFG) Nash equilibrium. This is the subject of the
companion paper \cite{hanzo-mean-field-alignment} and is treated here as a black box that, given a
population, returns an equilibrated one.

\emph{Middle (medium) --- the thermostat.} A scalar \emph{order parameter} $m \in [0,1]$ measures how
concentrated $p_t$ is (how aligned the fleet is). A coupling $J$ (how strongly each agent heeds the mean) and
a temperature $T$ (how much independent noise each agent injects) set $m$ through a ferromagnetic
self-consistency; the companion paper shows the useful operating point is \emph{critical}, $T \lesssim T_c$.

\emph{Outer (slow) --- the search.} A selection--recombination--mutation operator moves $p_t$ toward higher
fitness across \emph{generations}. This is the genetic algorithm and is the subject of this paper.
\end{definition}

The rest of the paper studies the outer loop and its coupling to the other two. The organising fact is that
all three are transport of the \emph{same} $p_t$; they differ only in the vector field they impose on it.

\section{A genetic algorithm is a fitness-driven Fokker--Planck flow}
\label{sec:ga-fpk}

Let $f:\mathcal{S}\to\R$ be fitness and write $\phi(p)=\E_{x\sim p}[f(x)]$ for the population-mean fitness.
A generational GA applies, in order: (i) \emph{selection}, reweighting $p$ toward high $f$; (ii)
\emph{crossover}, recombining sampled pairs; (iii) \emph{mutation}, perturbing offspring. Take the
weak-selection, small-step continuum limit of one generation. Selection and mutation become, respectively, a
drift and a diffusion of $p_t$ --- a Fokker--Planck equation --- and selection's drift is exactly the
replicator field.

\begin{proposition}[Selection is replicator drift; mutation is diffusion]
\label{prop:replicator}
In the weak-selection continuum limit, fitness-proportional (or Boltzmann, or tournament) selection on a
population $x=(x_i)$ of type frequencies induces the replicator dynamics \cite{taylor1978,hofbauer1998}
\begin{equation}
\dot{x}_i \;=\; x_i\bigl(f_i(x) - \phi(x)\bigr), \qquad \phi(x)=\sum_j x_j f_j(x),
\label{eq:replicator}
\end{equation}
and mutation at rate $\mu$ adds a diffusion term $\mu\,\Delta x_i$ (in genotype space), giving a
replicator--mutator (Fokker--Planck) flow. For frequency-independent $f$, mean fitness is non-decreasing
along \eqref{eq:replicator}: $\dot{\phi} = \mathrm{Var}_x(f) \ge 0$.
\end{proposition}

\begin{proof}[Proof sketch]
Fitness-proportional selection maps $x_i \mapsto x_i f_i/\phi$; subtracting $x_i$ and taking the small-step
limit gives $\dot x_i = x_i(f_i-\phi)/\phi$, which is \eqref{eq:replicator} after the standard time
reparametrisation $dt \mapsto \phi\,dt$. Boltzmann and tournament selection agree to first order in selection
strength (they differ only in the effective inverse temperature multiplying $f$). For the variance identity,
$\dot\phi=\sum_i \dot x_i f_i = \sum_i x_i f_i (f_i-\phi) = \E[f^2]-\E[f]^2$ when $f$ does not depend on $x$;
frequency dependence adds the cross term $\sum_i x_i f_i'\!\cdot\!\dot x$, which is where the game enters.
\end{proof}

The cross term is not a nuisance; it is the bridge. When fitness is \emph{frequency-dependent} ---
$f_i=f_i(x)$, an agent's payoff depends on what the rest of the fleet is doing --- \eqref{eq:replicator} is
evolutionary game theory, and its rest points are the Nash equilibria of the underlying game. That is the
inner loop of Definition~\ref{def:loops}: \textbf{replicator selection with frequency-dependent fitness is the
mean-field coupling.} The GA outer loop and the MFG inner loop are the same equation read at two speeds ---
selection slow, best-response fast --- so composing them is legitimate, not a metaphor. Mutation supplies the
diffusion that keeps $p_t$ full-support so selection has variance to act on (Proposition~\ref{prop:replicator}
stalls at $\mathrm{Var}=0$); crossover is the operator neither \eqref{eq:replicator} nor the mean field
contains, and it is where the speed comes from.

\section{Crossover is the speedup: compositionality and the schema theorem}
\label{sec:crossover}

A \emph{schema} $H$ is a subset of $\mathcal{S}$ fixing some coordinates and leaving others free; its
\emph{order} $o(H)$ is the number fixed and its \emph{defining length} $\delta(H)$ the span between the
outermost fixed coordinates. Holland's schema theorem \cite{holland1975} bounds the next-generation count $m(H,t)$ of individuals
matching $H$ under proportional selection, one-point crossover rate $p_c$, and mutation rate $p_m$:
\begin{equation}
\E\bigl[m(H,t+1)\bigr] \;\ge\; m(H,t)\,\frac{f(H)}{\phi}
\left(1 - p_c\,\frac{\delta(H)}{\ell-1} - o(H)\,p_m\right),
\label{eq:schema}
\end{equation}
where $\ell$ is the genome length and $f(H)$ the mean fitness of matching individuals. Read
\eqref{eq:schema} as a statement about \emph{building blocks}: a schema that is short ($\delta$ small),
low-order ($o$ small), and above average ($f(H)>\phi$) is disrupted little by crossover and grows roughly
geometrically. Crossover's job is to take two such blocks --- carried by different parents --- and place them
in one child.

Averaging (the mean field) cannot do this: it returns a convex combination of parents, which for a
combinatorial optimum is typically \emph{worse} than either. Independent search cannot do it either: it has
no channel to move a good block from one lineage to another. Crossover is the unique operator among the three
that transports building blocks \emph{between} lineages, and on a decomposable landscape that is decisive.

\begin{proposition}[Separable speedup]
\label{prop:speedup}
Let the fitness decompose additively over $M$ disjoint blocks, $f(x)=\sum_{b=1}^{M} g_b(x_b)$, each block of
size $k$ with a unique optimal setting. Suppose selection preserves any block already at its optimum in some
individual (each block-optimum is an above-average, short, low-order schema). Then uniform crossover assembles
the global optimum in $O(M\ln M)$ expected recombination events (coupon-collector over blocks), i.e.\ time
polynomial in $M$. A mutation-only search with no inter-block fitness gradient (the ``royal road'' class)
requires expected time exponential in $M$ to first reach the joint optimum.
\end{proposition}

\begin{proof}[Proof sketch]
Once each block's optimal setting is present in \emph{some} member of the population (which selection
maintains, since flipping a block off its optimum is selected against), a recombination event copies a parent's
optimal block into a child with constant probability per block. Collecting all $M$ optimal blocks into one
individual is the coupon-collector problem with $M$ coupons: expected $M H_M = O(M\ln M)$ draws, each a
recombination. For mutation-only on a flat-between-blocks (royal-road) landscape, no partial credit is
visible: the search is a random walk on $\{0,1\}^{Mk}$ whose only fitness signal is the all-optimal point, so
the hitting time is $\Theta(2^{\Theta(Mk)})$ up to the fraction of the space selection can prune, which the
royal-road construction makes vanishing. Hence polynomial versus exponential.
\end{proof}

\begin{remark}[The honest boundary: deception]
\label{rem:deception}
Proposition~\ref{prop:speedup} assumes block optima are \emph{above average in isolation}. A \emph{deceptive}
fitness \cite{goldberg1989} is built so that the block schemata leading toward the local optimum out-compete those
leading toward the global one; there crossover assembles the wrong blocks faster, and the GA converges
confidently to the trap. The correct fallback is mutation-and-diversity (raise $T$, Section~\ref{sec:order}),
not more crossover. The reason the framework nonetheless earns its keep for us is empirical: the fleet's
problems --- ``add these three actions to that tool,'' ``wire this module into the registry,'' ``port these
nine tools'' --- are \emph{separable by construction} (one file, one module, one tool per unit of work), which
is the regime where Proposition~\ref{prop:speedup}, not Remark~\ref{rem:deception}, governs.
\end{remark}

\section{The order parameter: premature convergence is the ordered phase}
\label{sec:order}

A GA's classic failure is \emph{premature convergence}: the population collapses onto one genotype, variance
vanishes, and by Proposition~\ref{prop:replicator} selection stalls --- the search freezes in a local optimum.
This is not a separate pathology from the companion paper's groupthink; it is the same event. Identify the
GA's population homogeneity with the ferromagnetic order parameter $m$: $m\to 1$ is a fully converged
population and a fully aligned fleet, one and the same ordered phase. Selection pressure is inverse
temperature; strong selection is cold and quenches, weak selection is hot and drifts. The cure is therefore
the cure for a quench: \emph{anneal}. Run hot early --- high mutation, weak coupling, diverse variants, so
crossover has distinct building blocks to find --- and cool toward criticality as fitness climbs, never all
the way to $m=1$. Simulated annealing \cite{kirkpatrick1983} is exactly this temperature schedule, and the companion paper's
critical operating point $T\lesssim T_c$ is where the annealed GA should spend its late generations:
correlated enough to converge, hot enough that a single high-fitness variant (or a single refuting critic) can
still tip the population out of a bad basin. The external field --- the directive --- breaks the symmetry so
the descent aims at the intended optimum rather than a merely-agreed one. \textbf{The GA is the search, the
ferromagnetic thermostat is how you keep the search from freezing, and they share one knob: $m$.}

\section{Embedding in enso training: the rollout group is a population}
\label{sec:enso}

The framework is not new machinery for the training stack; it is a name for a loop the stack already runs.
Group Relative Policy Optimization (GRPO) \cite{shao2024grpo}, implemented in the trainer of
\cite{hanzo-native-training} (\code{ml/hanzo-training/src/grpo/}), samples, for each
prompt, a \emph{group} of $G$ rollouts, scores them, and forms the group-relative advantage
$A_i=(r_i-\avg{r})/\mathrm{std}(r)$ before a length-normalised policy-gradient step with an optional $k3$ KL
penalty to a frozen reference. Read through Sections~\ref{sec:ga-fpk}--\ref{sec:order}:

\begin{itemize}[leftmargin=1.4em]
\item \textbf{Population $=$ the rollout group.} The $G$ samples are the generation; the policy is the
distribution $p_t$ they are drawn from.
\item \textbf{Selection $=$ the group-relative advantage.} $A_i$ is $f_i-\avg{f}$ with a variance
normaliser --- the replicator drift of \eqref{eq:replicator} verbatim: rollouts above the group mean are
reinforced, below are suppressed, and the group mean is the mean-field baseline. GRPO is replicator selection
with a frequency-defined baseline.
\item \textbf{Temperature $=$ sampling temperature; coupling $=$ the KL coefficient.} The $k3$ KL to the
reference is the coupling that keeps the group from quenching away from a sane prior; the sampling temperature
is $T$. Group size $G$ sets how well $\avg{r}$ estimates the mean field.
\item \textbf{Crossover $=$ expert recombination.} The AnyMoE in-engine gate trainer of
\cite{hanzo-native-training} (\code{amoe.rs}) trains gates over \emph{frozen quantized experts}: the experts
are reusable building blocks --- and, as the companion \emph{Streaming Mixture-of-Experts}
\cite{hanzo-streaming-experts} shows, independently addressable, byte-exact banks, exactly the substrate a
module-level crossover needs --- and the gate learns which to compose per token. Recombining experts across a
population of gates is crossover on the MoE genome; the island model of \cite{hanzo-intrinsic-islands} ---
independent sub-populations with occasional migration --- is the coarse-grained GA that keeps that
recombination from collapsing to one basin.
\item \textbf{Generations $=$ expert iteration.} Distilling the group's winning rollouts back into
\code{enso} (expert iteration / STaR \cite{zelikman2022star}) is the generational update --- the fittest
phenotypes of generation $t$ become the genotype prior of generation $t{+}1$; the harness of
Section~\ref{sec:harness} supplies the fitness that scores them.
\end{itemize}

\noindent\emph{The prescription that follows.} Because premature convergence is the ordered phase, the
training knobs $(G,\,T,\,\text{KL})$ should be \emph{annealed}, not fixed: start with large $G$, high $T$, low
KL (a hot, diverse group that finds building blocks), and cool $T$ and raise KL as reward climbs, holding the
group near criticality rather than driving the advantage to a spike. This is the training-time reading of
Section~\ref{sec:order}, and it is a config change over the existing GRPO loop, not new kernels.

\section{Embedding in the agentic harness: variants, critics, and grafts}
\label{sec:harness}

Inference-time, the fleet upgrades from \emph{align-and-best-respond} to
\emph{align-and-best-respond-and-evolve}. The mapping is exact and already partially in place in the workflow
engine that runs Hanzo's swarms:

\begin{center}
\begin{tabular}{@{}ll@{}}
\toprule
\textbf{GA operator} & \textbf{Harness realisation} \\
\midrule
population & $N$ solution \emph{variants}, not merely $N$ agents on one attempt \\
fitness $f$ & adversarial critics $+$ a green build $+$ the shipped-live outcome \\
selection & keep the variants the critics do not refute \\
\textbf{crossover} & a \emph{synthesis} agent: split candidates into modules, graft the fittest of each \\
mutation & a fresh-lens agent perturbing one module to escape a local optimum \\
generations & the loop-until-dry rounds; the completeness gap is the fitness gap \\
temperature & anneal coupling/selection each round: diverse early, aligned late \\
\bottomrule
\end{tabular}
\end{center}

Crucially this fitness is largely \emph{ground truth}, not a learned reward model that can be gamed: the
build either compiles or it does not, the tests either pass or they do not, and the ultimate signal is
whether the change ships and survives in production. That is what makes the harness landscape verifiable and
its separable component real, the low-epistasis regime Proposition~\ref{prop:speedup} needs.

Note the economy: \emph{the same critic is three things at once} --- the ferromagnetic fluctuation that stops
groupthink (Section~\ref{sec:order}), the GA's selection operator (Section~\ref{sec:ga-fpk}), and, when its
verdict is a scalar reward, the training signal of Section~\ref{sec:enso}. One mechanism, three roles, which
is why the flywheel closes.

\paragraph{A worked instance.} A hard tool port --- a browser tool with a $\sim$90-action surface and a
persistent-session requirement --- is exactly a separable problem with two hard blocks: the session model and
the action coverage. Generation~1 with no crossover (one agent, one attempt) tends to nail one block and stub
the other. The GA move is to spawn three variant implementations, score each by \{compiles, matches the
reference behaviour, survives the critic\}, and then \emph{crossover}: graft the best session model from one
variant onto the best action-coverage from another. By Proposition~\ref{prop:speedup} this reaches a correct
tool in generations logarithmic in the block count, where one agent grinding a single lineage does not.

\paragraph{The three-layer controller.} Stacking the loops of Definition~\ref{def:loops} gives a single
orchestrator: a \textbf{GA outer} loop (variants, critic-fitness, synthesis-crossover, annealed rounds) over a
\textbf{ferromagnetic middle} loop (hold $m$ near criticality, aim with the directive field) over an
\textbf{MFG inner} loop (each variant's agents best-respond to the shared mean). One knob --- $m$ / population
diversity --- threads all three, and the whole craft is annealing it from hot-diverse to critical-aligned
without quenching to $m=1$, with the directive aiming the descent. The companion paper builds the middle
layer; this paper builds the outer.

\section{What is real, what needs building, what is out of scope}
\label{sec:status}

\paragraph{Real.} The GRPO group, its group-relative advantage, the AnyMoE expert-gate trainer, and the
island decomposition are shipped and tested (\cite{hanzo-native-training}, \cite{hanzo-intrinsic-islands}); the
harness already runs populations of agents with adversarial critics as gates and a loop-until-dry completeness
check. Read as this paper prescribes, that is a genetic algorithm with selection and a fitness function.

\paragraph{Needs building.} The two operators the harness under-uses are \emph{crossover} (a synthesis agent
that decomposes candidates and grafts modules, rather than tournament-selecting a single winner) and an
explicit \emph{temperature schedule} (annealing coupling, mutation, and selection pressure round over round
with $m$ measured each generation). Both are orchestration logic over the existing workflow engine, not model
or kernel work.

\paragraph{Out of scope.} We do not claim a literal energy function for a fleet, nor that every task is
separable (Remark~\ref{rem:deception}); deceptive or strongly non-separable objectives fall back to
mutation-and-diversity and lose the crossover speedup. We do not treat multi-objective Pareto fronts, novelty
search, or open-ended evolution here; the target is the decomposable engineering objective with a scalar
gate, which is the fleet's daily work.

\section{Conclusion}
\label{sec:conclusion}

A fleet run as independent search wastes its members' partial wins; run as pure consensus it only averages.
Run as a genetic algorithm it does the thing neither can: recombine. We placed selection, mutation, and
crossover on one distribution $p_t$, showed selection is replicator drift and that frequency-dependent
selection is the mean-field coupling of the companion game (Proposition~\ref{prop:replicator}), and isolated
crossover as the operator responsible for the speedup --- provably superlinear on the separable landscapes the
fleet actually faces (Proposition~\ref{prop:speedup}), with deception as the named boundary
(Remark~\ref{rem:deception}). Premature convergence turned out to be the companion paper's ordered phase, so
the two papers share one control variable and one prescription: anneal $m$ toward criticality, never quench,
and aim with the directive. The framework is not speculative: GRPO's rollout group is the population, its
group-relative advantage is selection, AnyMoE experts are the building blocks, and the harness's critics are
already the fitness. What remains is to name and add the two missing operators --- a crossover-by-synthesis
step and an annealing schedule --- and let the fleet that trains the model, and the model that runs the fleet,
climb the same landscape together.

\begin{thebibliography}{9}

\bibitem{hanzo-native-training}
Hanzo AI --- ML Systems.
\newblock \emph{Native Training in hanzo-ml/hanzo-engine: One Engine, One Quant Format, One Backend}.
\newblock Hanzo Papers, 2026. Companion in the native-stack series.

\bibitem{hanzo-intrinsic-islands}
Hanzo AI --- ML Systems.
\newblock \emph{Intrinsic Islands: An Oracle-Anchored Escape Hatch for a Portable Kernel}.
\newblock hanzo-kernel Paper 07.3, \emph{One Source, Every Backend} series, 2026.

\bibitem{hanzo-streaming-experts}
Hanzo AI --- ML Systems.
\newblock \emph{Streaming Mixture-of-Experts: Bit-Exact Low-Memory Inference by Demand-Paging Routed Experts
from NVMe}.
\newblock Hanzo Papers, 2026. (Experts as independently addressable, byte-exact banks.)

\bibitem{hanzo-mean-field-alignment}
Hanzo AI --- ML Systems.
\newblock \emph{Mean-Field Alignment: Ferromagnetic Order and Critical Control of Agentic Fleets}.
\newblock hanzo-agentic Paper 01, \emph{Continually Improving Fleets} series, 2026. Companion to this paper.

\bibitem{shao2024grpo}
Z.\ Shao, P.\ Wang, Q.\ Zhu, et al.
\newblock DeepSeekMath: Pushing the Limits of Mathematical Reasoning in Open Language Models.
\newblock \emph{arXiv:2402.03300}, 2024. (Group Relative Policy Optimization, GRPO.)

\bibitem{zelikman2022star}
E.\ Zelikman, Y.\ Wu, J.\ Mu, and N.\ D.\ Goodman.
\newblock STaR: Bootstrapping Reasoning with Reasoning.
\newblock \emph{NeurIPS}, 2022. (Distil winning rollouts back into the model.)

\bibitem{holland1975}
J.\ H.\ Holland.
\newblock \emph{Adaptation in Natural and Artificial Systems}.
\newblock University of Michigan Press, 1975. (Schema theorem; building-block hypothesis.)

\bibitem{goldberg1989}
D.\ E.\ Goldberg.
\newblock \emph{Genetic Algorithms in Search, Optimization, and Machine Learning}.
\newblock Addison-Wesley, 1989. (Deception; the royal-road and building-block analyses.)

\bibitem{taylor1978}
P.\ D.\ Taylor and L.\ B.\ Jonker.
\newblock Evolutionarily stable strategies and game dynamics.
\newblock \emph{Mathematical Biosciences}, 40(1--2):145--156, 1978. (The replicator equation.)

\bibitem{hofbauer1998}
J.\ Hofbauer and K.\ Sigmund.
\newblock \emph{Evolutionary Games and Population Dynamics}.
\newblock Cambridge University Press, 1998. (Replicator dynamics; frequency-dependent selection.)

\bibitem{kirkpatrick1983}
S.\ Kirkpatrick, C.\ D.\ Gelatt, and M.\ P.\ Vecchi.
\newblock Optimization by Simulated Annealing.
\newblock \emph{Science}, 220(4598):671--680, 1983. (The temperature schedule.)

\end{thebibliography}

\end{document}
